\examyear{2012}
\examera{平成24年}
\examfield{代数}
\examgroup{B}
\examnumber{1}
\examtitle{平成24年 B 第1問}
\examtag{体論}
\examtag{ガロア理論}
\examtag{有理関数体}

\(L\) を複素数体 \(\C\) 上の \(4\) 変数有理関数体
\(\C(X_1,X_2,X_3,X_4)\) とする。体 \(L\) の \(2\) つの \(\C\) 自己同型 \(\sigma,\tau\) を
\[
  \sigma(X_1)=-X_1,\qquad
  \sigma(X_i)=X_i\quad (i=2,3,4),
\]
\[
  \tau(X_i)=X_{i+1}\quad (i=1,2,3),\qquad
  \tau(X_4)=X_1
\]
により定め，\(L\) の部分体 \(K\) を
\[
  K=\{a\in L\mid \sigma(a)=a,\ \tau(a)=a\}
\]
で定義する。

\begin{enumerate}
  \item \(L\) の \(K\) 上の拡大次数 \([L:K]\) を求めよ。
  \item 拡大 \(L/K\) の最大部分アーベル拡大を求めよ。
  \item \(L\) に含まれる \(K\) の \(2\) 次拡大をすべて求め，適当な \(L\) の元 \(\alpha\) を用いて \(K(\alpha)\) の形で表せ。
\end{enumerate}